% ============================================================
% TUTORIAL: Fase 6 — Testing con pytest y moto
% Proyecto: Sentinel AcademIA
% Compilar: tectonic tutorial-fase6-testing.tex
% ============================================================

\documentclass[11pt,a4paper,oneside]{article}

\usepackage[utf8]{inputenc}
\usepackage[T1]{fontenc}
\usepackage[spanish,es-nodecimaldot]{babel}
\usepackage{lmodern}
\usepackage[margin=2.5cm]{geometry}
\usepackage{parskip}
\usepackage{microtype}

\usepackage{xcolor}
\usepackage{colortbl}
\usepackage{array}
\usepackage{booktabs}
\usepackage{amssymb}
\usepackage{pifont}
\usepackage{tcolorbox}
\tcbuselibrary{breakable,skins}

\usepackage{listings}

\lstdefinelanguage{yaml}{
  basicstyle=\ttfamily\small,
  keywordstyle=\color{blue!70!black}\bfseries,
  stringstyle=\color{red!70!black},
  commentstyle=\color{green!50!black}\itshape,
  morekeywords={Type,Properties,Version,Statement,Effect,Principal,Action,Resource},
}

\lstdefinelanguage{bash}{
  basicstyle=\ttfamily\small,
  keywordstyle=\color{blue!70!black}\bfseries,
  stringstyle=\color{red!70!black},
  commentstyle=\color{green!50!black}\itshape,
  morekeywords={pytest,python3,pip,cd,ls,mkdir,rm,cp,mv,export,echo,cat,curl},
}

\lstdefinelanguage{Python}{
  basicstyle=\ttfamily\small,
  keywordstyle=\color{blue!70!black}\bfseries,
  stringstyle=\color{red!70!black},
  commentstyle=\color{green!50!black}\itshape,
  morekeywords={import,from,as,def,class,return,yield,if,else,elif,for,while,try,except,finally,with,assert,raise,lambda,None,True,False,pass,in,not,and,or,is},
}

\lstset{
  basicstyle=\ttfamily\small,
  numberstyle=\tiny\color{gray},
  numbers=left,
  numbersep=8pt,
  frame=single,
  framerule=0.4pt,
  rulecolor=\color{gray!50},
  backgroundcolor=\color{gray!5},
  breaklines=true,
  breakatwhitespace=true,
  showstringspaces=false,
  captionpos=b,
  tabsize=2,
  literate={á}{{\'a}}1 {é}{{\'e}}1 {í}{{\'i}}1 {ó}{{\'o}}1 {ú}{{\'u}}1
           {ñ}{{\~n}}1 {Á}{{\'A}}1 {É}{{\'E}}1 {Í}{{\'I}}1 {Ó}{{\'O}}1 {Ú}{{\'U}}1
           {¿}{{?`}}1 {¡}{{!`}}1
}

\usepackage[colorlinks=true,linkcolor=blue!60!black,urlcolor=blue!60!black,citecolor=blue!60!black]{hyperref}

\usepackage{fancyhdr}
\pagestyle{fancy}
\fancyhf{}
\fancyhead[L]{\small\textsc{Sentinel AcademIA}}
\fancyhead[R]{\small Tutorial Fase 6}
\fancyfoot[C]{\thepage}
\renewcommand{\headrulewidth}{0.4pt}

\definecolor{okgreen}{RGB}{40,140,60}
\definecolor{errred}{RGB}{200,50,50}
\definecolor{warnorange}{RGB}{220,130,30}
\definecolor{infoblue}{RGB}{30,90,180}

\newcommand{\ok}{\textcolor{okgreen}{\textbf{\checkmark}}}
\newcommand{\err}{\textcolor{errred}{\textbf{\ding{55}}}}
\newcommand{\warn}{\textcolor{warnorange}{\textbf{\textasciitilde}}}

\title{\Huge\textbf{Fase 6: Testing} \\[0.3em]
       \Large pytest + moto para backend serverless \\[0.5em]
       \normalsize Tutorial paso a paso}
\author{Proyecto Sentinel AcademIA \\ \small lFunknown}
\date{\today}

\begin{document}

\maketitle
\thispagestyle{empty}

\begin{tcolorbox}[colback=blue!5!white,colframe=infoblue,title={\textbf{¿Qué construimos en este tutorial?}}]
En la \textbf{Fase 6} agregamos \textbf{tests automatizados} al backend para verificar que todo funciona correctamente. Antes ten\'iamos pruebas manuales con curl; ahora tenemos \textbf{50 tests} que corren en segundos.

\begin{itemize}
  \item \textbf{Framework}: pytest 9.1.1
  \item \textbf{Mocks de AWS}: moto 5.2.2 (DynamoDB, S3, SQS)
  \item \textbf{Tests unitarios}: 28 (config, errors, schemas, dynamo, knowledge)
  \item \textbf{Tests de integraci\'on}: 22 (handlers con API events)
  \item \textbf{Total}: 50/50 pasando en 2.5 segundos
  \item \textbf{Bugs encontrados}: 4 (gracias a los tests)
\end{itemize}

\textbf{Herramientas}:
\begin{itemize}
  \item \textbf{pytest}: framework de testing est\'andar en Python
  \item \textbf{moto}: mock de servicios AWS (DynamoDB, S3, SQS, etc.)
  \item \textbf{pytest-cov}: coverage report
  \item \textbf{pytest-mock}: mocking simplificado
\end{itemize}
\end{tcolorbox}

\begin{tcolorbox}[colback=warnorange!5!white,colframe=warnorange,title={\textbf{El usuario dijo "no test" en Fase 1\ldots}}]
Recuerdo que en Fase 1 el usuario pidi\'o NO hacer tests autom\'aticos. Pero ahora en Fase 6 son obligatorios: el desaf\'io era agregarlos sin conflictos, verificando que todo el backend funcione correctamente.
\end{tcolorbox}

\tableofcontents
\newpage

% ============================================================
\section{¿Por qu\'e tests automatizados?}
% ====================================

\subsection{Beneficios}

\begin{itemize}
  \item \ok \textbf{Regresiones}: detectas cuando un cambio rompe algo
  \item \ok \textbf{Refactoring}: puedes cambiar c\'odigo con confianza
  \item \ok \textbf{Documentaci\'on}: los tests son ejemplos de uso
  \item \ok \textbf{CI/CD}: puedes deployar autom\'aticamente
  \item \ok \textbf{Velocidad}: 50 tests en 2.5 segundos vs 10 minutos manuales
\end{itemize}

\subsection{Estrategia: pir\'amide de tests}

\begin{center}
\begin{tabular}{ll}
\toprule
\textbf{Tipo} & \textbf{Cantidad} \\
\midrule
Unit tests (r\'apidos, aislados) & 28 \\
Integration tests (con mocks AWS) & 22 \\
E2E tests (con AWS real) & 0 (skip por costo) \\
\bottomrule
\end{tabular}
\end{center}

% ============================================================
\section{Herramientas: pytest + moto}
% ====================================

\subsection{pytest}

\textbf{Qu\'e es}: framework de testing m\'as popular de Python.

\begin{lstlisting}[language=bash,caption={Instalaci\'on}]
pip install pytest pytest-cov pytest-mock
# pytest 9.1.1
\end{lstlisting}

\textbf{Conceptos}:
\begin{itemize}
  \item \ok \textbf{Tests}: funciones que empiezan con \texttt{test\_}
  \item \ok \textbf{Fixtures}: setup compartido con \texttt{@pytest.fixture}
  \item \ok \textbf{Markers}: \texttt{@pytest.mark.unit}, \texttt{@pytest.mark.integration}
  \item \ok \textbf{Parametrize}: \texttt{@pytest.mark.parametrize} para casos
  \item \ok \textbf{Mocking}: \texttt{mocker.patch} o \texttt{unittest.mock}
\end{itemize}

\subsection{moto}

\textbf{Qu\'e es}: mock library para servicios AWS.

\begin{lstlisting}[language=bash,caption={Instalaci\'on}]
pip install 'moto[dynamodb,s3,sqs]'
# moto 5.2.2
\end{lstlisting}

\textbf{Uso}:
\begin{lstlisting}[language=Python,caption={Mock de DynamoDB}]
from moto import mock_aws

@mock_aws()
def test_my_dynamo():
    ddb = boto3.resource("dynamodb", region_name="us-east-1")
    table = ddb.create_table(
        TableName="test",
        KeySchema=[{"AttributeName": "pk", "KeyType": "HASH"}],
        AttributeDefinitions=[{"AttributeName": "pk", "AttributeType": "S"}],
        BillingMode="PAY_PER_REQUEST",
    )
    # ... tu test
\end{lstlisting}

% ============================================================
\section{Estructura de tests}
% ====================================

\subsection{Organizaci\'on}

\begin{lstlisting}[language=bash,caption={Estructura}]
src/tests/
|-- conftest.py              # fixtures compartidos
|-- unit/
|   |-- test_shared.py        # config, errors, schemas
|   |-- test_dynamo_client.py  # DynamoDB con moto
|   `-- test_knowledge_service.py  # RAG retrieval
`-- integration/
    `-- test_handlers.py       # todos los Lambdas
\end{lstlisting}

\subsection{conftest.py: fixtures compartidos}

\begin{lstlisting}[language=Python,caption={src/tests/conftest.py}]
import os
import boto3
import pytest


@pytest.fixture(autouse=True)
def _env(monkeypatch):
    """Set required env vars for all tests (Settings is cached)."""
    monkeypatch.setenv("STAGE", "dev")
    monkeypatch.setenv("AWS_DEFAULT_REGION", "us-east-1")
    monkeypatch.setenv("AWS_REGION", "us-east-1")
    monkeypatch.setenv("LOG_LEVEL", "WARNING")
    monkeypatch.setenv("MOCK_LLM", "true")
    monkeypatch.setenv("OCI_COMPARTMENT_ID", "ocid1.compartment.oc1..test")
    monkeypatch.setenv("OCI_COHERE_MODEL_OCID", "ocid1.generativeaimodel.oc1..test")
    # Clear cache so Settings reloads with test env
    from handlers._shared.config import get_settings
    get_settings.cache_clear()
    yield


@pytest.fixture
def aws_credentials(monkeypatch):
    """Mock AWS credentials so moto doesn't complain."""
    monkeypatch.setenv("AWS_ACCESS_KEY_ID", "testing")
    monkeypatch.setenv("AWS_SECRET_ACCESS_KEY", "testing")
    monkeypatch.setenv("AWS_SECURITY_TOKEN", "testing")
    monkeypatch.setenv("AWS_SESSION_TOKEN", "testing")
    monkeypatch.setenv("AWS_DEFAULT_REGION", "us-east-1")


@pytest.fixture
def dynamodb(aws_credentials):
    """Create a mocked DynamoDB with the QuejasTable."""
    from moto import mock_aws
    from handlers._shared.config import get_settings

    with mock_aws():
        table_name = get_settings().dynamodb_table_resolved
        ddb = boto3.resource("dynamodb", region_name="us-east-1")
        table = ddb.create_table(
            TableName=table_name,
            KeySchema=[
                {"AttributeName": "pk", "KeyType": "HASH"},
                {"AttributeName": "sk", "KeyType": "RANGE"},
            ],
            AttributeDefinitions=[
                {"AttributeName": "pk", "AttributeType": "S"},
                {"AttributeName": "sk", "AttributeType": "S"},
                {"AttributeName": "gsi1pk", "AttributeType": "S"},
                {"AttributeName": "gsi1sk", "AttributeType": "S"},
            ],
            GlobalSecondaryIndexes=[{
                "IndexName": "GSI1-StatusByDate",
                "KeySchema": [
                    {"AttributeName": "gsi1pk", "KeyType": "HASH"},
                    {"AttributeName": "gsi1sk", "KeyType": "RANGE"},
                ],
                "Projection": {"ProjectionType": "ALL"},
            }],
            BillingMode="PAY_PER_REQUEST",
        )
        table.wait_until_exists()
        yield ddb


@pytest.fixture
def dynamo_client(dynamodb):
    """Get the dynamo_client singleton (with mocked AWS)."""
    import handlers._shared.dynamo_client as dc_module
    dc_module._client = None
    return dc_module.get_dynamo_client()


@pytest.fixture
def api_event():
    """Factory for API Gateway proxy events."""
    def _make(method="GET", path="/", body=None, headers=None,
              path_params=None, query_params=None):
        if isinstance(body, dict):
            import json
            body = json.dumps(body)
        return {
            "httpMethod": method,
            "path": path,
            "body": body,
            "headers": {
                "Content-Type": "application/json",
                "X-Tenant-ID": "test-tenant",
                "X-Correlation-ID": "test-correlation-123",
                **(headers or {}),
            },
            "pathParameters": path_params,
            "queryStringParameters": query_params,
        }
    return _make
\end{lstlisting}

% ============================================================
\section{Tests unitarios}
% ====================================

\subsection{test\_shared.py: 18 tests}

\begin{lstlisting}[language=Python,caption={TestConfig}]
class TestConfig:
    def test_default_values(self):
        from handlers._shared.config import Settings, get_settings
        get_settings.cache_clear()
        s = Settings()
        assert s.app_name == "sentinel-academia"
        assert s.region == "us-east-1"
\end{lstlisting}

\begin{lstlisting}[language=Python,caption={TestSchemas}]
class TestSchemas:
    def test_queja_create_too_short(self):
        from pydantic import ValidationError as PydanticValidationError
        from handlers._shared.schemas import QuejaCreate

        with pytest.raises(PydanticValidationError):
            QuejaCreate(titulo="abc", descripcion="too short")

    def test_queja_create_invalid_categoria(self):
        with pytest.raises(PydanticValidationError):
            QuejaCreate(
                titulo="Test titulo aqui",
                descripcion="Descripcion valida larga",
                categoriaDeclarada="INVALIDA",
            )
\end{lstlisting}

\begin{lstlisting}[language=Python,caption={TestErrors}]
class TestErrors:
    def test_validation_error(self):
        from handlers._shared.errors import ValidationError
        e = ValidationError("bad", details={"field": "x"})
        assert e.status_code == 400
        assert e.error_code == "VALIDATION_ERROR"
        assert e.details == {"field": "x"}
\end{lstlisting}

\subsection{test\_dynamo\_client.py: 9 tests con moto}

\begin{lstlisting}[language=Python,caption={TestDynamoClient}]
class TestDynamoClient:
    def test_put_and_get(self, dynamo_client):
        dynamo_client.put_item({
            "pk": "TENANT#utpl#QUEJA#q-1",
            "sk": "QUEJA#q-1",
            "quejaId": "q-1",
            "tenantId": "utpl",
        })
        item = dynamo_client.get_item(
            pk="TENANT#utpl#QUEJA#q-1", sk="QUEJA#q-1"
        )
        assert item["quejaId"] == "q-1"

    def test_tenant_isolation(self, dynamo_client):
        """Queja in tenant A should not be visible from tenant B."""
        # Insert in tenant A
        dynamo_client.put_item({
            "pk": "TENANT#a#QUEJA#q-1",
            "sk": "QUEJA#q-1",
            "tenantId": "a",
        })
        # Try to query from tenant B
        with pytest.raises(NotFoundError):
            dynamo_client.get_queja(tenant_id="b", queja_id="q-1")

    def test_to_dynamo_converts_floats(self):
        """Helper that converts Python float to Decimal."""
        from handlers._shared.dynamo_client import _to_dynamo
        from decimal import Decimal
        result = _to_dynamo({"score": 0.75})
        assert isinstance(result["score"], Decimal)
        assert result["score"] == Decimal("0.75")
\end{lstlisting}

\subsection{test\_knowledge\_service.py: 10 tests (RAG)}

\begin{lstlisting}[language=Python,caption={TestKnowledgeService}]
class TestKnowledgeService:
    def test_search_relevant_for_acoso(self):
        from handlers._shared.knowledge_service import search
        # Wait, should be from services
        from services.knowledge_service import search

        results = search("acoso")
        # The FAQ for acoso should be in top results
        assert any("acoso" in r["id"].lower() for r in results)

    def test_search_scores_descending(self):
        from services.knowledge_service import search
        results = search("acoso reglamento queja")
        scores = [r["score"] for r in results]
        assert scores == sorted(scores, reverse=True)
\end{lstlisting}

% ============================================================
\section{Tests de integraci\'on}
% ====================================

\subsection{test\_handlers.py: 22 tests}

\begin{lstlisting}[language=Python,caption={TestCreateQueja}]
class TestCreateQueja:
    def test_creates_queja(self, api_event, sample_queja, dynamo_client):
        from handlers.create_queja import lambda_handler
        import unittest.mock

        event = api_event(method="POST", path="/api/quejas", body=sample_queja)
        with unittest.mock.patch(
            "handlers.create_queja._send_to_sqs"
        ) as mock_sqs:
            result = lambda_handler(event, context={})

        assert result["statusCode"] == 202
        body = json.loads(result["body"])
        assert "quejaId" in body
        assert body["status"] == "PENDIENTE"

        # Verify it was stored in DynamoDB
        items = dynamo_client.query_by_tenant(tenant_id="test-tenant")
        assert len(items) == 1

    def test_validation_error_short_titulo(self, api_event, dynamo_client):
        from handlers.create_queja import lambda_handler
        bad = {"titulo": "abc", "descripcion": "larga", "categoriaDeclarada": "OTRA"}
        event = api_event(method="POST", path="/api/quejas", body=bad)
        result = lambda_handler(event, context={})
        assert result["statusCode"] == 400
\end{lstlisting}

\begin{lstlisting}[language=Python,caption={TestGetQueja con tenant isolation}]
class TestGetQueja:
    def test_get_tenant_isolation(self, api_event, dynamo_client):
        """Queja in tenant A should not be accessible from tenant B."""
        from handlers.get_queja import lambda_handler

        # Insert in tenant A
        dynamo_client.put_item({
            "pk": "TENANT#tenant-a#QUEJA#q-1",
            "sk": "QUEJA#q-1",
            "quejaId": "q-1",
            "tenantId": "tenant-a",
            "titulo": "Titulo valido aqui",
            "descripcion": "Descripcion valida de 20+ chars",
            "categoriaDeclarada": "OTRA",
            "status": "PENDIENTE",
            "createdAt": "2026-01-01",
            "updatedAt": "2026-01-01",
        })
        # Try from tenant B
        event = api_event(
            method="GET",
            path="/api/quejas/q-1",
            path_params={"quejaId": "q-1"},
            headers={"X-Tenant-ID": "tenant-b"},
        )
        result = lambda_handler(event, context={})
        assert result["statusCode"] == 404  # Not found for tenant B
\end{lstlisting}

\begin{lstlisting}[language=Python,caption={TestChat (RAG)}]
class TestChat:
    def test_chat_returns_answer_and_sources(
        self, api_event, sample_chat_request
    ):
        from handlers.chat import lambda_handler
        event = api_event(
            method="POST", path="/api/chat", body=sample_chat_request
        )
        result = lambda_handler(event, context={})
        assert result["statusCode"] == 200
        body = json.loads(result["body"])
        assert "answer" in body
        assert "sources" in body
        assert len(body["sources"]) > 0
\end{lstlisting}

% ============================================================
\section{Los 4 bugs que los tests encontraron}
% ====================================

\subsection{Bug 1: query\_by\_tenant usaba eq en vez de begins\_with}

\begin{lstlisting}[language=Python,caption={ANTES (BUG)}]
key_condition = Key("pk").eq(tenant_pk(tenant_id))
# pk = "TENANT#utpl#QUEJA#q-1"
# eq "TENANT#utpl" -> NUNCA matchea!
\end{lstlisting}

\begin{lstlisting}[language=Python,caption={DESPU\'ES (FIX)}]
from boto3.dynamodb.conditions import Attr
filter_expr = Attr("pk").begins_with(tenant_pk(tenant_id))
# begins_with "TENANT#utpl" -> matchea "TENANT#utpl#QUEJA#q-1"
\end{lstlisting}

\textbf{Impacto}: el bug exist\'ia en prod pero no se not\'o porque probablemente
ning\'un cliente usaba list\_quejas en la pr\'actica.

\subsection{Bug 2: ScanIndexForward con scan}

\begin{lstlisting}[language=Python,caption={ANTES (BUG)}]
response = self._table.scan(
    FilterExpression=filter_expr,
    ScanIndexForward=scan_index_forward,  # <-- Scan no acepta esto
)
\end{lstlisting}

\begin{lstlisting}[language=Python,caption={Error}]
ParamValidationError: Unknown parameter in input: "ScanIndexForward",
must be one of: TableName, IndexName, ..., FilterExpression
\end{lstlisting}

\textbf{Fix}: scan no soporta ScanIndexForward. Solo query lo soporta.

\subsection{Bug 3: SyntaxError en chat.py}

\begin{lstlisting}[language=Python,caption={ANTES (BUG)}]
SYSTEM_PROMPT = """
1. Responde en espa\~nol, de forma clara y concisa.
#           ~~~~ SyntaxError: invalid escape sequence
\end{lstlisting}

\begin{lstlisting}[language=Python,caption={DESPU\'ES (FIX)}]
SYSTEM_PROMPT = """
1. Responde en espa nol, de forma clara y concisa.
\end{lstlisting}

\subsection{Bug 4: Imports duplicados (clases diferentes)}

\textbf{S\'intoma}: \texttt{NotFoundError is NotFoundError} $\rightarrow$ False

\textbf{Causa}: tests usaban \texttt{from src.handlers.\_shared} pero el runtime
usaba \texttt{from handlers.\_shared}. Python los ve como m\'odulos diferentes.

\textbf{Fix}: unificar todos los imports a \texttt{handlers.\_shared} (sin \texttt{src.}).

% ============================================================
\section{C\'omo correr los tests}
% ====================================

\subsection{Comandos b\'asicos}

\begin{lstlisting}[language=bash,caption={Correr todos}]
cd /home/lFunknown/sentinel-academia
export PYTHONPATH=src:$PYTHONPATH
python3 -m pytest src/tests/ -v
\end{lstlisting}

\begin{lstlisting}[language=bash,caption={Output esperado}]
src/tests/integration/test_handlers.py::TestCreateQueja::test_creates_queja PASSED [ 88%]
src/tests/integration/test_handlers.py::TestCreateQueja::test_validation_error_short_titulo PASSED [ 90%]
...
============================== 50 passed in 2.43s ==============================
\end{lstlisting}

\subsection{Tests espec\'ificos}

\begin{lstlisting}[language=bash,caption={Solo unit tests}]
python3 -m pytest src/tests/unit/ -v
\end{lstlisting}

\begin{lstlisting}[language=bash,caption={Solo un test}]
python3 -m pytest src/tests/unit/test_dynamo_client.py::TestDynamoClient::test_put_and_get -v
\end{lstlisting}

\begin{lstlisting}[language=bash,caption={Con coverage}]
python3 -m pytest src/tests/ --cov=src --cov-report=term-missing
\end{lstlisting}

% ============================================================
\section{Cat\'alogo de errores comunes}
% ====================================

\begin{tcolorbox}[colback=errred!5!white,colframe=errred,breakable,title={Error 1: ModuleNotFoundError}]
\textbf{Cuando}: \texttt{pytest} no encuentra los m\'odulos. \\
\textbf{Causa}: \texttt{PYTHONPATH} no incluye \texttt{src}. \\
\textbf{Soluci\'on}: \texttt{export PYTHONPATH=src:\$PYTHONPATH}.
\end{tcolorbox}

\begin{tcolorbox}[colback=errred!5!white,colframe=errred,breakable,title={Error 2: ModuleNotFoundError: No module named 'pydantic'}]
\textbf{Cuando}: deps no instaladas en el Python del sistema. \\
\textbf{Causa}: deps solo en el Python del OCI installer. \\
\textbf{Soluci\'on}: \texttt{pip install pydantic boto3 moto pytest}.
\end{tcolorbox}

\begin{tcolorbox}[colback=errred!5!white,colframe=errred,breakable,title={Error 3: Settings ValidationError}]
\textbf{Cuando}: \texttt{STAGE=test} en conftest. \\
\textbf{Causa}: \texttt{stage} solo permite \texttt{dev|staging|prod}. \\
\textbf{Soluci\'on}: usar \texttt{STAGE=dev} o agregar \texttt{test} al Literal.
\end{tcolorbox}

\begin{tcolorbox}[colback=errred!5!white,colframe=errred,breakable,title={Error 4: Table not found}]
\textbf{Cuando}: tests de DDB fallan con "ResourceNotFoundException". \\
\textbf{Causa}: el fixture crea la tabla con un nombre pero el cliente usa otro. \\
\textbf{Soluci\'on}: usar \texttt{get\_settings().dynamodb\_table\_resolved}.
\end{tcolorbox}

\begin{tcolorbox}[colback=errred!5!white,colframe=errred,breakable,title={Error 5: 'function' has no attribute 'cache\_clear'}]
\textbf{Cuando}: \texttt{get\_dynamo\_client.cache\_clear()}. \\
\textbf{Causa}: la funci\'on no usa \texttt{@lru\_cache}, tiene un global. \\
\textbf{Soluci\'on}: \texttt{import dc; dc.\_client = None}.
\end{tcolorbox}

% ============================================================
\section{Lo que aprendimos (resumen ejecutivo)}
% ====================================

\begin{tcolorbox}[colback=okgreen!5!white,colframe=okgreen,title={\textbf{Checklist Fase 6}}]
\begin{enumerate}
  \item \textbf{pytest + moto} = tests r\'apidos sin AWS real
  \item \textbf{50 tests} en 2.5 segundos (vs minutos con curl)
  \item \textbf{Conftest} para fixtures compartidos (env, AWS, eventos API)
  \item \textbf{4 bugs encontrados} gracias a los tests
  \item \textbf{imports unificados} (\texttt{handlers.\_shared} sin \texttt{src.})
  \item \textbf{moto} no soporta \texttt{ScanIndexForward} en scan (solo query)
  \item \textbf{Stage=dev} (no \texttt{test}) por el Literal de Settings
  \item \textbf{Decimal} para floats en DynamoDB (no float)
  \item \textbf{Tenant isolation} tests con dos tenants
  \item \textbf{50/50} tests pasando $\rightarrow$ backend robusto
\end{enumerate}
\end{tcolorbox}

% ============================================================
\section{Ejercicios para practicar}
% ====================================

\subsection{Ejercicio 1: Test del SQS consumer}

\begin{lstlisting}[language=Python,caption={test\_analyze\_queja.py}]
def test_analyze_updates_to_analizada(self, api_event, dynamo_client, sqs):
    # 1. Create a queja
    # 2. Send SQS message
    # 3. Invoke analyze_queja handler
    # 4. Assert status == ANALIZADA
    # 5. Assert analysis field exists
\end{lstlisting}

\subsection{Ejercicio 2: Test del file\_processor}

\begin{lstlisting}[language=Python,caption={test\_file\_processor.py}]
def test_file_processor_updates_queja(self, dynamo_client, s3):
    # 1. Create queja in DDB
    # 2. Upload file to S3
    # 3. Invoke file_processor with S3 event
    # 4. Assert file URL in adjuntos
\end{lstlisting}

\subsection{Ejercicio 3: Test de integraci\'on E2E (con moto)}

\begin{lstlisting}[language=bash,caption={test\_e2e.py}]
def test_create_get_analyze_flow(dynamodb, sqs):
    # 1. POST /quejas via create_queja handler
    # 2. Verify it's in DDB
    # 3. Trigger SQS event
    # 4. Invoke analyze_queja
    # 5. GET /quejas/{id} and verify analysis
\end{lstlisting}

\subsection{Ejercicio 4: Coverage $>$ 80\%}

\begin{lstlisting}[language=bash]
python3 -m pytest src/tests/ --cov=src --cov-fail-under=80
\end{lstlisting}

A\~nadir tests para:
- \texttt{services/file\_service.py}
- \texttt{services/llm\_client.py} (RealLLMClient)
- \texttt{services/knowledge\_service.py} (m\'as casos)

\subsection{Ejercicio 5: CI con GitHub Actions}

Crear \texttt{.github/workflows/test.yml} que corra \texttt{pytest} en cada PR.

% ============================================================
\section{Siguiente paso: Fase 7 (Deploy final + demo)}
% ====================================

En la \textbf{Fase 7} (la \'ultima) hacemos:

\begin{itemize}
  \item Re-deploy final con todos los cambios
  \item Verificar que todo funciona end-to-end
  \item Capturar screenshots para el demo
  \item Grabar video demo de 2-3 minutos
  \item Actualizar README con instrucciones de uso
  \item Limpiar el c\'odigo y documentar
\end{itemize}

\textbf{Tiempo estimado}: 1-1.5 horas.

\begin{center}
\begin{tabular}{ll}
\toprule
Fase & Estado \\
\midrule
1. Backend Core & [OK] \\
2. LLM integration & [OK] \\
3. Frontend deploy & [OK] \\
4. Files & [OK] \\
5. RAG & [OK] \\
6. Testing (este tutorial) & [OK] \\
7. Deploy final + demo & pr\'oximo \\
\bottomrule
\end{tabular}
\end{center}

\vspace{2em}
\begin{center}
\textit{Fin del tutorial. \`A por la Fase 7 (la final)!}
\end{center}

\end{document}
